\section{Future Research Directions}
\label{sec:future-directions}

Future CV and ITS deployments should choose RSU spacing, antenna geometry, and
supplemental coverage from the fraction of compact J2735
messages completed within the application deadline. The evaluation must include urban intersections,
high-rise street canyons, highways, and sparsely served roads because distance
alone does not characterize obstruction or coverage.

At the matched large-object condition, vehicle-originated content required
substantially longer than the same-sized edge-returned content. CAV developers
should therefore select, compress, partition, or progressively transmit
vehicle-originated state according to the available uplink, the current radio
condition, and the decision deadline. These mechanisms also need
explicit expiration and recovery so that delayed fragments do not consume
resources after the application can no longer use the complete object.

For mobile network operators and 5G/6G researchers, CAV demand changes in
arrival time, dominant direction, and duration. A driving event can create an
immediate uplink- or downlink-heavy burst. Remote teleoperation can then sustain
video, state, and control traffic for the remainder of the assistance session.
Other deadline-sensitive operations may begin while this stream remains
active. The same-device measurements show that a nominal TDD ratio does not
determine end-to-end burst latency or sustained goodput. The load experiments
also show that a continuing stream can move compact foreground exchanges across
their deadlines. Future radio work should make balanced and uplink-oriented
frame structures interoperable across chipsets, vehicle gateways, and
infrastructure while preserving attachment and service state. Schedulers must
consider each workload's direction, deadline, and expected duration while
isolating simultaneous traffic. Common benchmarks should combine asynchronous
bursts with sustained streams across independently scheduled UEs while varying
signal condition, TDD allocation, mobility, and handover. They should report
direction-separated RTT, deadline completion over all issued requests, upload
and download application goodput, fairness, configuration-state provenance,
and recovery gaps
\cite{ngmn2021tdd,itu2023m2160}.
